\section{Introduction}\label{sec:security-intro}

Healthcare information systems are primary targets for security breaches due to the high commercial value of Protected Health Information (PHI). 
We addressed these risks by managing multi-layered security challenges, including web-based access, mobile authentication, sensitive patient consent delegation, and complex doctor-patient workflows. 
Furthermore, the system maintains audit trails and supports real-time session revocation.

Since traditional monolithic security models are insufficient for such a distributed system CFI-Care adopts a \textbf{Defense-in-Depth} strategy, leveraging industry standards such as \ac{OAuth} 2.0 and \ac{OIDC} with \ac{JWT}. 
Security is enforced through distributed validation across the system firewall and the \ac{API} gateway (Nginx), and the \ac{BFF} layer, ensuring multiple checkpoints before data access.

\section{Architecture Overview}\label{sec:security-architecture}

The security architecture follows a layered model where each component provides a specialized validation check. 
Figure~\ref{fig:security-components} illustrates the interaction between edge clients, the identity provider, and backend resources. 

\begin{figure}[H]
    \centering
    \includegraphics[width=\linewidth]{Security/images/components-ver.png}
    \caption{Security Architecture Components: Edge clients (Browser, Mobile), API Gateway (Nginx, OAuth2 proxy), Backend-for-Frontend (Node.js), Identity (Keycloak), Storage (Redis, Postgres), and Resource Server (HAPI FHIR).}
    \label{fig:security-components}
\end{figure}

The following pillars define the implementation:

\begin{itemize}
    \item \textbf{Identity \& Access Management:} Keycloak serves as the central Identity Provider (IdP), issuing signed \acp{JWT}. 
    Authentication pathways are structurally split by client type to optimize token security for both browser-based environments and native mobile deployments
    \item \textbf{Transport \& Edge Security:} Nginx terminates TLS and enforces HTTPS-only communication. 
    At the application layer, defensive browser headers and strict cookie policies are enforced at the BFF boundary to neutralize client-side vulnerability vectors.
    \item \textbf{Distributed Validation:} Token integrity, scope claims, and patient-specific bindings are validated at the BFF layer before requests reach the resource storage, ensuring that the principle of least privilege is maintained.
\end{itemize}

% Healthcare information systems are prime targets for security breaches due to
% the commercial value of protected health information (PHI) . CFI-Care platform
% presents multi-layered challenges: securing web-based browser access, mobile
% device authentication, sensitive patient consent grants, and doctor-patient
% access workflows.Also maintaining audit trails is required and so is enabling
% session revocation

% Traditional monolithic security approaches are insufficient for this use case
% Instead, we adopted a defense-in-depth strategy leveraging modern standards
% (\ac{OAuth} 2.0/OIDC with \ac{JWT}) and distributed validation across the
% \ac{API} gateway (nginx), backend-for-frontend (Node.js/BFF) and \ac{FHIR}
% server for final checks before database access

% The CFI-Care security architecture follows a defense-in-depth model with
% multiple validation layers Figure~\ref{fig:security-components} provides an
% overview of the security components and their interactions.Figures
% \ref{fig:browser-tokens-flow}, \ref{fig:mobile-tokens-flow}, and
% \ref{fig:token-flow-overview} illustrate the key authentication and
% authorization flows across browser, mobile workflows, this section details the
% design and implementation principles


% \begin{figure}[H]
%     \begin{center}
%         \includegraphics[width=\linewidth]{Security/images/components-ver.png}
%         \caption{Security Architecture Components: Edge clients (Browser, Mobile), API Gateway (Nginx, OAuth2 proxy), Backend-for-Frontend (Node.js), Identity (Keycloak), Storage(Redis, Postgres) and Resource Server (HAPI FHIR) with their security-focused interactions.}
%         \label{fig:security-components}
%     \end{center}
% \end{figure}

% \paragraph{Identity \& Access (Keycloak + Tokens)}
% OAuth 2.0/OIDC issues JWTs signed by Keycloak. Browser and mobile follow
% different integration paths: browser via gateway-managed OIDC
% (\texttt{oauth2-proxy} + nginx, Authorization Code + PKCE), mobile via
% Authorization Code + PKCE in-app. Tokens include user identity, scopes, roles,
% and claims.

% \paragraph{Transport \& Edge (TLS + Headers)}
% nginx terminates TLS, enforces HTTPS-only access, and sets forwarded headers.
% The BFF applies Helmet security headers (CSP, HSTS, X-Frame-Options)


